\documentclass{article}
\usepackage[english]{babel}
\usepackage{xcolor}

\title{Advanced Machine Learning\\Report 1}

\author{Mohammed Zaid Shaikh\\Jan Wyrzykowski\\Tomasz Żochowski}

\begin{document}

\maketitle
\tableofcontents

\section{Introduction}

\textcolor{red}{something about importance of dealing with unlabeled data}

\textcolor{blue}{idk it kinda sounds like an abstract}
Our goal was to evaluate different techniques related to filling missing labels in data. Firstly, we prepared datasets for logistic regression and created new target columns according to various missing data schemes. Then, we implememted the FISTA algorithm which we used in Logistic Lasso Regression. Finally, we used different label infering methods to reconstruct original labels and compared the quality of classification between these methods as well as the Naive approach and the Oracle approach.


\section{Optimization with FISTA}

\textcolor{red}{Found the link to the paper which introduces FISTA: https://www.ceremade.dauphine.fr/~carlier/FISTA}


\section{Labeling missing data}


\section{Results}


\end{document}
